\begin{remark}[E13: which of a scope's two connectives binds tighter]\label{rem:precedence}
Remark~\ref{rem:disjunction-detail} asserts that Datadog's \texttt{or} binds more loosely
than its comma, so \texttt{\{a:x or a:y,*\}} parses with the comma inside the right
branch. The justification offered was an analogy with the Lucene-style filter syntax
Datadog's editor accepts, which is not evidence about this grammar, and the claim decides
how $160$ \texttt{SCOPE\_BOOL} occurrences are \emph{translated} rather than whether they
are. No corpus query discriminates: the only shape mixing the two connectives puts
match-all terms on the other side, so both readings agree on it.

Three series over two tag keys --- \texttt{aa:x bb:2}, \texttt{aa:y bb:1},
\texttt{aa:x bb:1} --- and one question, \texttt{\{aa:x or aa:y, bb:1\}}. Under the
looser-\texttt{or} reading this is $aa{:}x \vee (aa{:}y \wedge bb{:}1)$ and selects the
first series whatever its \texttt{bb}; under the tighter reading it is
$(aa{:}x \vee aa{:}y) \wedge bb{:}1$ and excludes it. The account selected it. The
assertion was right, and is now a measurement.
\end{remark}

\begin{remark}[E14: what a negated match-all selects]\label{rem:negated-all}
Nine corpus occurrences carry \texttt{\{*,!*\}}, and \S\ref{sec:gaps} put them in the
unimplemented column on the strength of two claims: that PromQL can denote an empty
selection, and that \texttt{\{*,!*\}} \emph{is} one. The first is true --- a
contradictory matcher pair returns nothing, read off a real server. The second was never
asked.

It is false. On the account, both \texttt{\{*,!*\}} and \texttt{\{!*\}} select the
series. Whatever \texttt{!*} does to a scope, it is not exclusion, and the nine
occurrences are not a request for the empty set. They stay in the unimplemented column,
where they already were, but for a different and much duller reason: the term is
discarded. The entry in \S\ref{sec:gaps} that called them an empty selection was
reasoning about a construct nobody had run.

That reading was one observation short, and the shortfall is worth recording because it
is the same shape as the error it corrects. A series coming back is consistent with two
readings --- the term was understood and excluded nothing, or the term was never
evaluated --- and points alone cannot separate them. Every \texttt{/api/v1/query}
response names the expression Datadog evaluated, and in both cases that expression is
\texttt{sum:qlport.e14.negall\{run:\ldots\}}: the term is absent. It did not match
everything. It was dropped before evaluation.

Three controls make that a contrast rather than a lone observation, because
\emph{selects the series} says nothing until something has failed to. Over the same
metric and window, a value the series does not carry selects nothing, and a negation of
a value it does carry selects nothing --- so exclusion is expressible here. And
\texttt{!rig:zzz}, a negation of an absent value, selects the series \emph{and survives
into the echoed expression}. Negation works and is reported; this spelling of it is
neither.

This is the kind of error that costs nothing in any number the paper reports and would
have cost a reader everything: a translation of those nine queries built on the
catalogue's reading would have emitted an expression returning no data where Datadog
returns a series. The instrument that settles it was free, present in every response the
harness had already cached, and read for the first time in the audit that produced this
paragraph; it is now recorded for every probe the harness runs.
\end{remark}

\begin{remark}[E15: whether a \texttt{rate} metric passes through $\phi_\rho$ unchanged]\label{rem:rate-identity}
Definition~\ref{def:phi} sends every non-counter series through $\phi_\rho$ by the
identity, and Remark~\ref{rem:type-measured} counts $134$ corpus occurrences naming
metrics the vendor declares as \texttt{rate}. Those occurrences sit inside $\kappa$ and
their exactness rests on that clause, which was asserted: a Datadog \texttt{RATE} is
submitted as a per-second value over a declared interval, and whether what the system
stores and reduces is that value or the count it implies is exactly the question
$\phi_\rho$ exists to ask.

One \texttt{RATE} series and one \texttt{COUNT} series, both carrying five points of the
same constant at interval $60$, both read back with an \texttt{avg} and a \texttt{sum}
rollup over a window spanning them:

\begin{center}\footnotesize
\begin{tabular}{@{}lrr@{}}
\toprule
& \texttt{RATE} & \texttt{COUNT} \\
\midrule
\texttt{.rollup(avg, 300)} & $7.0$ & $7.0$ \\
\texttt{.rollup(sum, 300)} & $28.0$ & $28.0$ \\
\bottomrule
\end{tabular}
\end{center}

\noindent
Identical, and both the submitted values unrescaled --- the $28$ rather than $35$ being
the window catching four of the five points, which is the ordinary boundary effect and
not a factor. So a \texttt{RATE} is stored and reduced as the numbers it was given, the
identity clause is what the system does, and the $134$ occurrences rest on a
measurement.
\end{remark}

\begin{remark}[E16: what a keyless scope term globs over]\label{rem:keyless-glob}
This experiment was designed to settle a reading and instead exposed a defect in this
project's own corpus pipeline, which is why it is worth reporting even though the
question it asked turned out not to arise.

\S\ref{sec:regular} had read a keyless scope term such as \texttt{*-location} as a glob
over the whole \texttt{key:value} string, and that reading has an odd consequence:
\texttt{cluster-location:us-east1} does not end in \texttt{-location}, so the dashboards
carrying the term would select nothing. Three series were submitted, one per tag key ---
\texttt{cluster-location:eu}, \texttt{site-location:us}, \texttt{region:ap} --- so that
a glob over keys and a glob over whole tags would give different answers.

Neither answer came back. Datadog's query parser \emph{rejects} the keyless term
outright, both bare and with a value --- \texttt{\{*-location\}} and
\texttt{\{*-location:*\}} --- giving the same error it gives for a single-character
wildcard. A leading \texttt{*} in a keyless term is not Datadog
syntax, so no corpus query could have contained one --- and the $31$ that appeared to
were an artefact of \texttt{corpus.\_TMPL} truncating \texttt{\$cluster-location} at its
hyphen. The experiment asked what a construct means and learned that the construct does
not exist, which is a better outcome than either reading.
\end{remark}
